\documentclass[11pt,a4paper]{article}
\usepackage[utf8]{inputenc}
\usepackage[T1]{fontenc}
\usepackage[a4paper,margin=2.3cm]{geometry}
\usepackage{amsmath,amssymb}
\usepackage{xcolor}
\usepackage{hyperref}
\usepackage{enumitem}
\hypersetup{colorlinks=true,linkcolor=blue,urlcolor=blue}
\setlength{\parskip}{0.45em}\setlength{\parindent}{0pt}
\sloppy

\newcommand{\rev}[1]{{\color{black}#1}}
\newcommand{\resp}[1]{{\color{red}#1}}
\newcommand{\Frec}{F_\mathrm{rec}}
\newcommand{\Fest}{F_\mathrm{est}}

\begin{document}

\begin{center}
{\Large\textbf{Point to Point Response to the Reviewers}}\\[0.3em]
Manuscript: \emph{Blind Catalytic Quantum Error Correction: Target-State Estimation and Fidelity Recovery Without A Priori Knowledge}\\
Authors: Hikaru Wakaura, Taiki Tanimae \quad Journal: \emph{Quantum} \quad Manuscript ID: [to be filled]
\end{center}

\resp{We sincerely appreciate the editor and reviewers' valuable and insightful comments. We have highlighted portions of the manuscript that have been updated. Reviewer comments are reproduced in full in black; our responses follow in red.}

\resp{\textbf{Summary of the revision.} Both reviewers identified the same central issue: the numerical recovery step in the submitted version was the PSD projection of the estimate, so the reported estimation--recovery agreement was built into the implementation, and the manuscript described this as a simulated catalytic channel. We completely agree. Following the alternative that both reviewers explicitly offered, the revised manuscript is reframed as a \emph{classical target-estimation study under an explicitly stated idealized recovery interface}: the recovery map is now defined once, in Eq.~(1), as the mode-inclusion-restricted PSD projection of the estimate, the manuscript states in the Abstract, Sec.~2.2, Sec.~2.3, Sec.~5.3, and Sec.~6.4 that the catalytic covariant transformation itself is not simulated, and every fidelity is described as an upper bound for a physical implementation. All numerical results were regenerated from a single forward/inverse-consistent noise module; the fixed-point crossover equation, the ``single-copy'' resource claims, the tomographic-saturation claims, and the $d \geq 128$ data were removed; the catalytic theorem was restated; the Lipschitz lemma was replaced by the rigorous one-sided bound that Reviewer~2 kindly derived; and the repository was frozen with relative paths, raw results files, and a script-to-figure map. Details follow.}

\section*{Reviewer 1}

\rev{\textbf{Summary: what are the main questions posed by the manuscript and how does it answer them?} The manuscript asks whether catalytic quantum error correction (CQEC) can be used when the ideal target state is not known in advance. It proposes a two-stage method: estimate a target state from a noisy density matrix, then use that estimate as the target of a density-matrix-level CQEC/ICEC recovery step. It compares several estimators and reports a near-linear relation between estimation and recovery fidelity, together with dimension- and copy-dependent performance claims.}

\resp{We thank the reviewer for this accurate summary.}

\rev{\textbf{Assessment.} I do not recommend acceptance in the present form. I would like to raise three main technical concerns.}

\rev{\textbf{First,} the implemented recovery does not match the physical description. In \texttt{blind\_cqec/recovery.py}, \texttt{icec\_recover(noisy, target\_estimate)} ignores \texttt{noisy} and returns \texttt{psd\_project(target\_estimate)}, which clips negative eigenvalues and renormalizes the trace. The public estimators already return PSD-normalized states, so this recovery step normally leaves the estimate unchanged. Consequently, in the packaged \texttt{icec\_recover} path, equality between estimation and recovery fidelity is built into the implementation. Other benchmark scripts call \texttt{ICECProtocol} from \texttt{icec.py}. In the \texttt{catalytic\_transform} method called by \texttt{ICECProtocol.correct}, however, the output is produced by matrix-level amplification and interpolation toward the supplied target. This benchmark path does not construct a joint system-catalyst output state or specify an explicit CPTP covariant map whose system and catalyst marginals are verified, and it does not numerically verify covariance. Theoretical catalytic reachability may justify an idealized recovery assumption, but it does not make this interpolation a numerical implementation or validation of the physical channel described in the manuscript.}

\rev{The manuscript and the documentation in \texttt{blind\_cqec/recovery.py} state that full rank makes mode inclusion automatic. This is not correct as stated. Full rank does not imply that coherent modes are present: the maximally mixed state is full rank but has no off-diagonal coherence. Thus full rank alone does not guarantee mode inclusion.}

\resp{We thank the reviewer for these comments. We completely agree on both points.
(i) The recovery step in the submitted version was indeed the PSD projection of the estimate, and describing it as a simulated catalytic transformation was not appropriate. We have adopted the reviewer's suggested alternative (item~1 of the suggested changes): the manuscript now states, at the first point where recovery is introduced, that the recovery map is the density-matrix-level \emph{interface}
$\rho_\mathrm{rec} = \Pi(\mathcal{M}_{\rho_\mathrm{noisy}}(\hat\rho_\mathrm{est}))$ [Eq.~(1), Sec.~2.3], where $\mathcal{M}$ removes estimated coherences on modes absent from the noisy state and $\Pi$ is the PSD projection; that the catalytic covariant transformation, catalyst, joint state, and covariance are \emph{not} simulated (Sec.~2.2, new text); and that consequently all reported fidelities are upper bounds and the estimation--recovery agreement of Sec.~5.3 is a consistency check of the interface rather than an empirical discovery (Sec.~5.3, rewritten; Fig.~4 caption). The packaged \texttt{icec\_recover} now uses \texttt{noisy} to enforce the mode-inclusion constraint (package v0.2.0, \texttt{CHANGELOG.md}), and the legacy \texttt{ICECProtocol} path used by the algorithm benchmarks is driven by that same restricted target, so the two code paths no longer implement different models. The Limitations section (Sec.~6.4, items 1--2) lists the explicit CPTP catalytic map as the first open problem.
(ii) The full-rank statement was wrong, and we have removed it everywhere. Theorem~1 is restated (Sec.~2.1) with the explicit remark that the mode set is determined by which off-diagonal elements are nonzero and is unrelated to rank, with the full-rank diagonal state as the counterexample; Sec.~3.2 now explains that coherence maximization generally \emph{introduces} modes absent from the noisy state, that depolarizing noise does not create missing off-diagonal elements, and that such modes are removed by $\mathcal{M}$ before recovery; the package docstring no longer contains the claim. Mode survival is checked explicitly against the threshold $10^{-10}$ stated in Sec.~7.}

\rev{\textbf{Second,} the analytical bound concerns the projection map, not the catalytic map described earlier. For a valid density matrix, this projection leaves the target estimate unchanged. Thus the bound establishes output-state stability for the PSD post-processing surrogate, but it does not establish continuity or finite-resource stability for an explicitly constructed catalytic recovery channel.}

\resp{We thank the reviewer for this comment and completely agree. The bound is now presented (Sec.~2.4) as a statement about the interface $\Pi\circ\mathcal{M}$ only: Eq.~(3) is the direct Fuchs--van de Graaf bound for density-matrix estimates (which is the case for every estimator in the paper), and Lemma~1 / Eq.~(4) is the one-sided projection bound for general Hermitian inputs. The paragraph following Eq.~(4) states explicitly that these bounds ``do not establish continuity or finite-resource stability of a physically constructed, target-dependent catalytic channel; that is a separate open problem,'' and the Abstract and Introduction no longer describe the result as a Lipschitz property of catalytic recovery.}

\rev{\textbf{Third,} the claimed tomographic and sample-complexity conclusions are not supported by an operational sampling model. \texttt{benchmark\_blind\_qem\_compare.py} averages repeated evaluations of the same deterministic noisy density matrix. In \texttt{benchmark\_blind\_sensitivity.py}, the added matrix perturbation is itself scaled by \texttt{1/sqrt(n\_copies)} before \texttt{n\_copies} samples are averaged. In addition, the alternative \texttt{ICECProtocol} implementation explicitly inserts copy-number dependence through a \texttt{sqrt(n\_copies)} coherence-amplification rule and an interpolation weight that increases with \texttt{n\_copies}. Scaling extracted from these constructions is therefore at least partly prescribed by the simulation rather than obtained from a measurement-and-reconstruction experiment. It cannot establish saturation of a tomographic bound or a superlinear physical copy advantage.}

\resp{We thank the reviewer for this careful reading and completely agree. The revised manuscript makes no claim of saturating a tomographic bound or of a physical superlinear copy advantage. Sec.~5.11 is retitled ``Sample-complexity context'' and now states that the copy counts characterize only the fluctuation-averaging step of our perturbation model, conditional on oracle density-matrix access, and that an operational shot-level study is required before any such claim can be made; Fig.~15 is relabeled with its reference lines described as visual guides with illustrative prefactors. Sec.~5.4 discloses both mechanisms the reviewer identified---the fixed-amplitude Gaussian perturbation ensemble and the $n$-dependence contributed by the recovery convergence model---and restricts interpretation of the fitted exponents to the within-condition ordering of strategies, with each exponent labeled by its exact noise condition. The interpretive figure that previously accompanied these exponents was removed. Building a Born-rule measurement model is listed as the second open problem in Sec.~6.4.}

\rev{\textbf{To what extent have you checked the technical correctness of the submission?} My assessment is based on a targeted, rather than exhaustive, examination of the manuscript source, its tests, and the benchmark scripts most directly related to the concerns above. Although the repository also contains separate ICEC prototype tests, I did not find validation of an explicit joint CPTP covariant recovery map used for the reported benchmarks. I did not verify every proof, audit every source file, or regenerate every figure and table. Several paper-level benchmark scripts appear to require path and import cleanup before they can be run from a fresh checkout.}

\resp{We thank the reviewer. We confirm that no explicit joint CPTP covariant map is validated in the benchmarks, and the manuscript now says so directly. The path and import issues are addressed below.}

\rev{\textbf{Presentation.} The manuscript is readable, but it should distinguish access to a noisy quantum system, the measurements needed to infer its density matrix, classical target estimation, and physical recovery. The claimed ``single-copy'' overhead is not operationally established when the estimator takes a complete noisy density matrix as input, unless the measurements, repeated preparations, and reconstruction cost needed to obtain that matrix are included. The mode-inclusion assumption should also be stated more carefully. Full rank does not imply that all coherent modes are present, so it does not by itself make mode inclusion automatic. The bibliography should also be audited; the active entry for arXiv:2603.11568 contains a concrete author/title mismatch.}

\resp{We thank the reviewer for these comments and agree with all of them. (i) The four layers are now separated explicitly: oracle density-matrix access (Abstract; Sec.~3.2; Sec.~5.10 scope statement; Table~7 caption; Sec.~6.4 item~2), classical estimation (Sec.~3), the idealized recovery interface (Sec.~2.2--2.3), and the unsimulated physical recovery (Sec.~6.4 item~1). (ii) Every ``single-copy'' claim has been removed from the Abstract, Table~7, and Conclusion and replaced by ``oracle access to the noisy density matrix, whose acquisition costs $\Theta(d^2/\epsilon^2)$ shots''; Sec.~5.11 additionally includes the contractivity argument (also raised by Reviewer~2) showing that universal single-copy recovery is impossible for any deterministic target-independent apparatus. (iii) The mode-inclusion statement is corrected as described above. (iv) The reviewer is correct; this was an error. The entry for arXiv:2603.11568 has been corrected to J.~Raghoonanan and T.~Byrnes, ``Hybrid Quantum Error Correction and Mitigation by Purification,'' and the corresponding text in Sec.~6.3 now describes that work (purification quantum error suppression, PQES) accurately.}

\rev{\textbf{Reproducibility.} In the supplied repository, I did not find complete reproduction instructions mapping the supplied scripts to every submitted figure and table. Several paper-level scripts contain hard-coded local paths and use different recovery models. I also did not find separate raw numerical-data files corresponding to the submitted figures, although the data-availability statement says that these data are publicly available.}

\resp{We thank the reviewer for pointing this out and have corrected all three issues in the frozen repository accompanying this resubmission: every hard-coded author-machine path has been replaced by paths relative to the script location, so all scripts run from a fresh checkout; every benchmark script writes its numerical results to \texttt{blind\_cqec\_pkg/results/*.json}, which are shipped in the repository; and \texttt{REPRODUCE.md} maps every figure and table of the manuscript to its generating script, results file, and runtime. All scripts now import the same noise module and the same mode-restricted recovery target, so there is a single recovery model (Sec.~7, ``Unified noise module''). Undeclared dependencies are declared under the \texttt{bench} extra in \texttt{pyproject.toml}.}

\rev{\textbf{Suggested changes.} 1.~Identify the recovery map used for every result \ldots 2.~Separate estimator performance from recovery performance \ldots 3.~Limit the Lipschitz claim to the PSD post-processing surrogate \ldots 4.~Rebuild the tomographic and sample-complexity study from an explicit measurement and estimation model \ldots 5.~Correct the manuscript and package-documentation claim that full rank makes mode inclusion automatic. Correct the bibliographic entry for arXiv:2603.11568, unify the code paths, remove local paths, and release the raw data and reproduction commands. The estimator comparison may form the basis of a useful revised study, but the present evidence does not numerically demonstrate an explicit catalytic covariant recovery channel or establish the claimed finite-copy and single-copy operational performance of blind CQEC.}

\resp{We thank the reviewer for this constructive list. Items 1, 3, and 5 are implemented as described above. Item~2: Sec.~5.3 now reports $\Fest$ and $\Frec$ separately for all 84 conditions and localizes their difference to the mode-restriction penalty (largest deviation: $\Fest = 1.00$ vs.\ $\Frec = 0.41$ for the Regev state under dephasing-only noise); the effect of PSD projection alone is negligible for our estimators, which already output density matrices. Item~4: we have not built a measurement-and-reconstruction model in this revision; instead we removed every conclusion that depended on one (tomographic saturation, superlinear copy advantage, single-copy overhead) and state the shot-level study as required future work (Secs.~5.11, 6.4). We agree with the reviewer's closing assessment and have reframed the contribution accordingly as an estimator-design study under an explicitly assumed ideal recovery interface.}

\section*{Reviewer 2}

\rev{\textbf{Summary.} The manuscript proposes making target-dependent catalytic quantum error correction (CQEC) ``blind'' by estimating the ideal state from a noisy output and using that estimate as the recovery target. It compares five estimators, studies estimation--recovery fidelity correlations and noise/dimension dependence, and claims low-copy state recovery with advantages over quantum error mitigation. Its evidence comprises analytical bounds and density-matrix benchmarks, including algorithmic outputs, a VQE example, and a circuit-simulator check.}

\resp{We thank the reviewer for this accurate summary.}

\rev{\textbf{Assessment.} I recommend rejection in the present form. Estimating a target before catalytic recovery is a reasonable research direction, but this manuscript does not establish the advertised physical protocol or resource advantages. The main deficiencies concern the input information and the recovery operation. The Lipschitz proof has separate mathematical problems, although the one-sided bound needed for Eq.~(7) can be repaired.}

\rev{\textbf{1.} The numerical recovery does not implement a catalytic transformation. Section 2.4 defines $\omega = \Pi(\hat\sigma)$, where $\Pi(A) = A_+/\mathrm{Tr}(A_+)$ clips negative eigenvalues and renormalizes. The recovery function ignores the noisy input and returns the projected estimate. The principal packaged estimators already output density matrices, so $\omega = \hat\sigma$ up to roundoff. Their estimation--recovery fidelity agreement is consequently an identity. No catalyst, joint state, covariant CPTP channel, or catalyst reuse is simulated. The legacy code instead interpolates toward the specified target and copies the catalyst unchanged. An estimator comparison assuming ideal recovery could be legitimate, but cannot validate that assumed recovery or establish that estimation is the ``sole bottleneck.''}

\resp{We thank the reviewer for this comment and completely agree. As described in our response to Reviewer~1 (first concern), the manuscript is now explicitly an estimator comparison under an assumed ideal recovery interface, defined once in Eq.~(1) and declared as unsimulated physics in Secs.~2.2, 2.3, 5.3, and 6.4. The phrase ``sole bottleneck'' has been removed; Sec.~5.3 and the Conclusion now say that, \emph{within this interface}, performance decomposes into an estimation problem plus the mode-survival ceiling $\Delta_\mathrm{mode}$, and that whether a physical catalytic implementation preserves this structure is open. The recovery function now enforces mode inclusion using the noisy state, so $\Frec \leq \Fest$ with a data-dependent gap (Fig.~4: $r = 0.999$ at $d = 4$ falling to $r = 0.770$ at $d = 64$), which the text presents as a verification of Eq.~(3) rather than as a discovery.}

\rev{\textbf{2.} Complete matrix access is not single-copy quantum access. All estimators receive the full noisy density matrix as classical data. One physical copy does not reveal its populations and phases. Obtaining them experimentally requires a measurement model and a sample budget. Classical nonlinear estimation is legitimate, but cannot be identified with a deterministic quantum channel acting on one unknown state; formal channel inverses are generally not CPTP. A simple limitation illustrates why the distinction matters. Let $\sigma_\pm = (I \pm X)/2$. Dephasing produces $\rho_\pm = (I \pm \eta X)/2$, with $0 < \eta < 1$. Both retain nonzero coherence. Any fixed deterministic apparatus using one noisy copy and an initially independent, target-independent catalyst induces a CPTP map $R$. Trace-distance contractivity gives $D(R(\rho_+), R(\rho_-)) \leq \eta$, while $D(\sigma_+, \sigma_-) = 1$. Thus at least one recovery error is at least $(1 - \eta)/2$. This excludes arbitrarily accurate universal recovery under those single-copy assumptions. It does not exclude many-copy estimation, additional target information, improvements on other restricted ensembles, or conditional protocols with stated success probabilities. If ``single-copy'' describes only an online stage after training, all training costs must be included.}

\resp{We greatly appreciate the reviewer providing this argument. We completely agree, and we have incorporated the contractivity limitation verbatim in spirit into Sec.~5.11 (with attribution in this letter): it now states that universal single-copy recovery to arbitrary accuracy is impossible for any deterministic target-independent apparatus, that the estimators exploit exactly the additional information supplied by oracle density-matrix access, and that all resource claims must include that cost. Every ``single-copy'' statement has been removed (Abstract, Table~7, Sec.~5.10, Conclusion), and Sec.~3.2 now states explicitly that a density matrix is not a single measurement outcome. We have also removed the identification of channel inversion with a quantum channel; it is described as classical post-processing throughout.}

\rev{\textbf{3.} The catalytic existence theorem is misstated. Shiraishi--Takagi, Theorem 2 gives arbitrarily accurate correlated-catalytic conversion under a sufficient resonant-mode inclusion condition, and exact conversion when the target is full rank. It does not give the manuscript's exact ``if and only if'' statement with a full-rank input. Moreover, existence for each specified source--target pair does not provide a target-independent recovery apparatus or its resource cost. Full rank also does not imply all coherence modes: a full-rank diagonal state is a counterexample. Depolarization does not create missing off-diagonal elements. The estimator may therefore introduce modes for which the cited sufficient condition has not been established.}

\resp{The reviewer is correct; this was an error. Theorem~1 (Sec.~2.1) is restated as the correlated-catalytic conversion statement of Shiraishi--Takagi Theorem~2: mode inclusion is sufficient for conversion to arbitrary accuracy, with exactness when the \emph{target} is full rank; the ``if and only if'' and the full-rank-input condition are removed. The text following the theorem now makes both of the reviewer's further points explicitly: that the theorem is an existence statement per source--target pair and supplies neither a target-independent apparatus nor a cost, and that the mode set is unrelated to rank. Sec.~3.2 states that coherence maximization generally introduces modes absent from the noisy state, that depolarization does not create missing off-diagonal elements, and that such modes are removed by $\mathcal{M}$ so that the effective target satisfies mode inclusion by construction.}

\rev{\textbf{4.} The Lipschitz criticism must be distinguished from the recoverable fidelity bound. The proof uses $\|A_+ - B_+\|_1 \leq \|A - B\|_1$, which is false for noncommuting Hermitian matrices. [Explicit $3\times3$ trace-one counterexample given, with $\|A_+ - B_+\|_1 \approx 0.02227682 > 0.02119824 \approx \|A - B\|_1$.] Nevertheless, for Hermitian trace-one $A$ and any density matrix $\sigma$, the needed one-sided bound is valid: $\|\Pi(A) - \sigma\|_1 \leq 2\|A - \sigma\|_1$. Indeed, write $A = A_+ - A_-$ and $s = \mathrm{Tr}(A_-)$. Then $\|\Pi(A) - A\|_1 = 2s$, whereas $\|A - \sigma\|_1 \geq 2s$, by testing the negative spectral projector and using $\mathrm{Tr}(A - \sigma) = 0$. The triangle inequality proves the bound. Combined with $1 - F(\omega,\sigma) \leq \|\omega - \sigma\|_1$ for squared fidelity, this recovers Eq.~(7) within the classical projection model. Consequently, Eq.~(7) should not be described as disproved. Its repair does not establish stability of an actual target-dependent catalytic channel. Fixed-channel contractivity cannot supply that missing result, and a fidelity-regression slope is not a trace-norm Lipschitz constant. The stated general derivation of Eq.~(1) also incorrectly substitutes oracle infidelity for its trace-norm error; a valid argument requires additional justification.}

\resp{We greatly appreciate the reviewer pointing out both the counterexample and the valid one-sided bound. We verified the counterexample numerically ($\|A_+ - B_+\|_1 = 0.022277 > 0.021198 = \|A - B\|_1$) and have adopted the reviewer's argument as Lemma~1 in Sec.~2.4, with the proof exactly as the reviewer sketched (orthogonal supports of $A_\pm$ giving $\|\Pi(A) - A\|_1 = 2s$; the negative spectral projector giving $\|A - \sigma\|_1 \geq 2s$). The two-sided clipping inequality has been removed, and the text states that it is false for noncommuting inputs. The corollary is now Eq.~(4), $\Frec \geq 1 - 2\|\hat\rho_\mathrm{est} - \rho_\mathrm{target}\|_1 - 2\Delta_\mathrm{mode}$, which reduces to Eq.~(3) with constant~1 for density-matrix inputs. The paragraph after Eq.~(4) limits the interpretation to the interface and states that a regression slope is not a Lipschitz constant. The former Eq.~(1), which substituted oracle infidelity for its trace-norm error, has been deleted; the fidelity bound is now derived directly from Fuchs--van de Graaf for the interface and does not reference an oracle map.}

\rev{\textbf{5.} The claimed analytical crossover fails its own algebra. Equation (8) states $d^\star = 1/[2(\gamma_{AD} + \gamma/d^\star)]$. With $\gamma = 1$ and $\gamma_{AD} = 0.1$, this gives $0.2 d^\star + 2 = 1$, hence $d^\star = -5$, not 25--40. The preceding conversion from elementwise estimates to trace-norm scaling is also unjustified. This derivation cannot support the headline crossover claim.}

\resp{The reviewer is correct; this was an error. The fixed-point equation and the value $d^\star \approx 25$--$40$ have been removed from the Abstract, Introduction, Sec.~2.5, Sec.~5.9, and Conclusion. Sec.~2.5 now gives only the qualitative scaling argument for why coherence maximization degrades with dimension and why exact channel inversion does not, and states that the location of any crossover is empirical and depends on calibration accuracy. With the unified noise implementation, the regenerated data (Sec.~5.6) show no crossover at all under exact calibration---channel inversion coincides with the oracle at every dimension---and a crossover appears only under parameter misspecification (Sec.~5.7); the manuscript now reports exactly that.}

\rev{\textbf{6.} The benchmarks do not validate the statistical or operational claims. Repeated copies of a fixed state after a deterministic channel have the same density matrix; averaging those matrices does not produce finite-sample improvement. The scripts instead use ad hoc matrix perturbations or change target ensembles across copy counts. Reruns give exponents near 0.17 and 0.12, consistent with Fig.~6's plotted labels but inconsistent with its superlinear claims. The packaged inverse undoes the full combined channel and can recover the target exactly in deterministic arithmetic, whereas legacy benchmarks invert only selected components and use different noise models. The reported estimator crossover may therefore reflect inconsistent implementations. The QEM comparison substitutes matrix manipulations for operational PEC, ZNE, and virtual distillation, without matched sampling costs. The circuit check does not implement catalytic recovery; one simulator shot still supplies an exact density matrix, and its fitted noise parameter uses the ideal target. The VQE example likewise uses classical matrix postprocessing and exact expectation values. These demonstrations do not establish physical state recovery or the claimed resource advantage.}

\resp{We thank the reviewer for this thorough analysis and agree on every point. (i) The reviewer's diagnosis of the inconsistent implementations was exactly right: the legacy dimension-sweep script applied uniform dephasing while inverting with the gap-dependent factor, and did not invert amplitude damping. All scripts now import a single noise module, forward and inverse (Sec.~7; unit tests enforce element-wise agreement and round-trip exactness), and every affected result was regenerated. The previously reported crossover disappeared as the reviewer anticipated. (ii) The superlinear-scaling and tomographic-saturation claims are withdrawn; Sec.~5.4 discloses the perturbation model and the convergence-model $n$-dependence and restricts interpretation to within-condition orderings (see also Reviewer~1, third concern). (iii) Sec.~5.10 opens with a scope statement that the QEM comparison is a non-resource-matched density-matrix idealization of each method, reports that idealized PEC reaches $F = 1.000$ at every dimension, and makes no deployment-level claim. (iv) Sec.~5.13 is retitled ``Circuit-derived states: an estimator exercise with oracle calibration,'' discloses that $p_\mathrm{eff}$ is calibrated from the ideal state and that no recovery circuit is executed, and no longer describes the exercise as a validation. (v) Sec.~5.12 describes the H$_2$ example as classical post-processing with exact expectation values, a proof of principle, and additionally reports the LiH null result.}

\rev{\textbf{Technical correctness.} I read the complete manuscript and cover letter, checked the foundational theorem, the crossover algebra, and the stability arguments. I verified the counterexamples, derived the valid one-sided bound, and checked the single-copy limitation under its stated assumptions.}

\resp{We are grateful for the depth of this examination; the corrected theorem, the one-sided bound, and the contractivity limitation are all now in the manuscript.}

\rev{\textbf{Presentation.} The motivation, estimator taxonomy, and broad organization are understandable. However, the presentation repeatedly conflates an existence theorem, classical density-matrix estimation, and an implementable quantum channel. This prevents a broader quantum-information reader from reconstructing the claimed protocol. The ``single-copy,'' ``threshold-free,'' and ``sole bottleneck'' claims are insufficiently qualified. The text also assigns joint-system $O(d^6)$ cost to calculations whose released recovery is an $O(d^3)$ matrix projection, and describes circuit-level confirmation without a recovery circuit. The related-work discussion should better address state learning, inverse-channel estimation, correlated catalysis, and sampling limits. These are substantive clarity problems, beyond stylistic editing.}

\resp{We thank the reviewer for these comments and agree. The three layers are now kept separate throughout (see our response to Reviewer~1, presentation). ``Single-copy'' and ``sole bottleneck'' are removed; ``threshold-free'' is retained only as a property of the reachability theorem, with the operational qualifications stated in Sec.~6.1. Sec.~2.2 and Sec.~7 now state that the released recovery is an $O(d^3)$ projection and that the $O(d^6)$ joint-system simulation is not performed. The circuit exercise is no longer described as a confirmation. Regarding related work, the revised Sec.~3.2 and Sec.~5.11 discuss the relation to quantum state estimation and sampling limits, and Sec.~2.1 addresses correlated catalysis; we recognize that a fuller treatment of state learning and inverse-channel estimation would strengthen the paper and would welcome the reviewer's guidance on specific references.}

\rev{\textbf{Reproducibility.} Only partially sufficient. The code is public and executable, and some reported values are reproducible. However, the package and legacy scripts implement different algorithms and noise models. Several scripts depend on author-machine paths and undeclared dependencies, and the promised raw data and generated figure artifacts were absent from the repository.}

\resp{We thank the reviewer and have addressed each point as described in our response to Reviewer~1 (reproducibility): unified noise model and recovery target across package and scripts, relative paths, declared dependencies, shipped \texttt{results/*.json}, and a script-to-figure map in \texttt{REPRODUCE.md}.}

\rev{\textbf{Suggested changes.} A scientifically meaningful resubmission would require substantial new work: define the target ensemble, physical inputs, measurements, side information, output, success probability, and total estimation/recovery costs; construct a covariant CPTP recovery on the joint system and catalyst, with mode conditions, implementation error, catalyst preparation, and reuse costs---alternatively, restrict the paper to classical estimation under an explicitly assumed ideal recovery oracle; correct the catalytic theorem and replace the general Lipschitz lemma with the valid one-sided bound, clearly limiting its interpretation; establish physical-channel stability separately and repair Eq.~(1); remove or rederive the crossover; rebuild copy-scaling experiments from measurement records and compare methods under matched operational budgets; release a frozen implementation with consistent noise models, dependencies, data, and figure-generation instructions. My recommendation rests primarily on the missing evidence for the advertised physical protocol and resource advantages. Repairing the projection bound alone would not resolve those deficiencies.}

\resp{We thank the reviewer for this clear roadmap. Of the two alternatives offered, we have taken the second: the revised manuscript restricts itself to classical estimation under an explicitly assumed ideal recovery interface, and no longer advertises a physical protocol or resource advantage. The theorem is corrected, the one-sided bound replaces the lemma with its interpretation limited, the former Eq.~(1) is removed, the crossover is removed, the copy-scaling and QEM claims that required measurement records or matched budgets are withdrawn and stated as future work, and the implementation is frozen with consistent models, dependencies, data, and figure instructions. We agree that constructing and validating the joint covariant CPTP recovery, and rebuilding the resource analysis from measurement records, are the substantive next steps; they are stated as such in Sec.~6.4 and the Conclusion, and are beyond the scope of what the present, deliberately narrowed, manuscript claims.}

\end{document}
